\section{卷五 107-150}

\subsection{第107篇： 五种要向神稱讚的人 1-43}
\subsubsection{稱讚神将在曠野漂流的人領到可居住的城邑 1-9}
你們要稱謝耶和華，因他本為善；他的慈愛永遠長存！
願耶和華的贖民說這話，就是他從敵人手中所救贖的，
從各地，從東從西，從南從北，所招聚來的。
他們在曠野荒地漂流，尋不見可住的城邑，
又飢又渴，心裡發昏。
於是，他們在苦難中哀求耶和華；他從他們的禍患中搭救他們，
又領他們行走直路，使他們往可居住的城邑。
但願人因耶和華的慈愛和他向人所行的奇事都稱讚他；
因他使心裡渴慕的人得以知足，使心裡飢餓的人得飽美物。

\subsubsection{稱讚神折断死蔭的捆鎖，将人領出黑暗 10-16}
那些坐在黑暗中、死蔭裡的人被困苦和鐵鍊捆鎖，
是因他們違背神的話語，藐視至高者的旨意。
所以，他用勞苦治服他們的心；他們仆倒，無人扶助。
於是，他們在苦難中哀求耶和華；他從他們的禍患中拯救他們。
他從黑暗中和死蔭裡領他們出來，折斷他們的綁索。
但願人因耶和華的慈愛和他向人所行的奇事都稱讚他；
因為他打破了銅門，砍斷了鐵閂。

\subsubsection{稱讚神医治愚妄人因過犯和罪孽的疾病苦楚，脱离死亡 17-22}
愚妄人因自己的過犯和自己的罪孽便受苦楚。
他們心裡厭惡各樣的食物，就臨近死門。
於是，他們在苦難中哀求耶和華；他從他們的禍患中拯救他們。
他發命醫治他們，救他們脫離死亡。
但願人因耶和華的慈愛和他向人所行的奇事都稱讚他。
願他們以感謝為祭獻給他，歡呼述說他的作為！

\subsubsection{稱讚神救人脱离海上的狂风，领人到所去的海口 23-31}
在海上坐船，在大水中經理事務的，
他們看見耶和華的作為，並他在深水中的奇事。
因他一吩咐，狂風就起來，海中的波浪也揚起。
他們上到天空，下到海底；他們的心因患難便消化。
他們搖搖幌幌，東倒西歪，好像醉酒的人；他們的智慧無法可施。
於是，他們在苦難中哀求耶和華，他從他們的禍患中領出他們來。
他使狂風止息，波浪就平靜。
風息浪靜，他們便歡喜；他就引他們到所願去的海口。
但願人因耶和華的慈愛和他向人所行的奇事都稱讚他。
願他們在民的會中尊崇他，在長老的位上讚美他！

\subsubsection{稱讚神使曠野乾旱之地成为可居住的有福之地，使窮乏人生养众多 32-43}
他使江河變為曠野，叫水泉變為乾渴之地，
使肥地變為鹼地；這都因其間居民的罪惡。
他使曠野變為水潭，叫旱地變為水泉。
他使飢餓的人住在那裡，好建造可住的城邑，
又種田地，栽葡萄園，得享所出的土產。
他又賜福給他們，叫他們生養眾多，也不叫他們的牲畜減少。
他們又因暴虐、患難、愁苦，就減少且卑下。
他使君王蒙羞被辱，使他們在荒廢無路之地漂流。
他卻將窮乏人安置在高處，脫離苦難，使他的家屬多如羊群。
正直人看見就歡喜；罪孽之輩必塞口無言。
凡有智慧的，必在這些事上留心，也必思想耶和華的慈愛。

\subsection{第108篇(大衛)： 清早起来赞美神,求神帮助百姓攻击敌人 1-13}
\subsubsection{清早起来赞美神的慈爱和诚实 1-5}
神啊，我心堅定；我口（榮耀）要唱詩歌頌！
琴瑟啊，你們當醒起！我自己要極早醒起！
耶和華啊，我要在萬民中稱謝你，在列邦中歌頌你！
因為，你的慈愛大過諸天；你的誠實達到穹蒼。
神啊，願你崇高過於諸天！願你的榮耀高過全地！
\subsubsection{求神起来帮助百姓攻擊敵人 6-13}
\begin{table}[H]
\centering
\begin{tabular}{p{8.5cm}|p{8.75cm}}\hline\hline
\multicolumn{1}{c|}{神的應允}&\multicolumn{1}{c}{向主祈求}\\\hline
&求你應允我們，用右手拯救我們，好叫你所親愛的人得救。\\\hline
神已經指著他的聖潔說（應許我）：我要歡樂；我要分開示劍，丈量疏割谷。
基列是我的；瑪拿西是我的；以法蓮是護衛我頭的；猶大是我的杖；
摩押是我的沐浴盆；我要向以東拋鞋；我必因勝非利士呼喊。
&誰能領我進堅固城？誰能引我到以東地？
神啊，你不是丟棄了我們嗎？神啊，你不和我們的軍兵同去嗎？
求你幫助我們攻擊敵人，因為人的幫助是枉然的。
我們倚靠神才得施展大能，因為踐踏我們敵人的就是他。\\\hline
\end{tabular}
\end{table} 


\subsection{第109篇(大衛-基督的被弃):因仇敵無故辱罵攻擊，而向神禱告 1-31}
\subsubsection{被仇敵用撒謊和詭詐的舌頭攻擊依然專心祈禱 1-5}
我所讚美的神啊，求你不要閉口不言。
因為惡人的嘴和詭詐人的口已經張開攻擊我；他們用撒謊的舌頭對我說話。
他們圍繞我，說怨恨的話，又無故地攻打我。
他們與我為敵以報我愛，但我專心祈禱。
他們向我以惡報善，以恨報愛。
\subsubsection{愿神報應攻擊自己的仇敵 6-20}
願你派一個惡人轄制他，派一個對頭站在他右邊！
他受審判的時候，願他出來擔當罪名！願他的祈禱反成為罪！
願他的年日短少！願別人得他的職分！
願他的兒女為孤兒，他的妻子為寡婦！
願他的兒女漂流討飯，從他們荒涼之處出來求食！
願強暴的債主牢籠他一切所有的！願外人搶他勞碌得來的！
願無人向他延綿施恩！願無人可憐他的孤兒！
願他的後人斷絕，名字被塗抹，不傳於下代！
願他祖宗的罪孽被耶和華記念！願他母親的罪過不被塗抹！
願這些罪常在耶和華面前，使他的名號斷絕於世！
因為他不想施恩，卻逼迫困苦窮乏的和傷心的人，要把他們治死。
他愛咒罵，咒罵就臨到他；他不喜愛福樂，福樂就與他遠離！
他拿咒罵當衣服穿上；這咒罵就如水進他裡面，像油入他的骨頭。
願這咒罵當他遮身的衣服，當他常束的腰帶！
這就是我對頭和用惡言議論我的人從耶和華那裡所受的報應。
\subsubsection{求神恩待搭救，我要用口極力稱謝讚美耶和華 21-31}
主─耶和華啊，求你為你的名恩待我；因你的慈愛美好，求你搭救我！
因為我困苦窮乏，內心受傷。
我如日影漸漸偏斜而去；我如蝗蟲被抖出來。
我因禁食，膝骨軟弱；我身上的肉也漸漸瘦了。
我受他們的羞辱，他們看見我便搖頭。
耶和華─我的神啊，求你幫助我，照你的慈愛拯救我，
使他們知道這是你的手，是你─耶和華所行的事。
任憑他們咒罵，惟願你賜福；他們幾時起來就必蒙羞，你的僕人卻要歡喜。
願我的對頭披戴羞辱！願他們以自己的羞愧為外袍遮身！
我要用口極力稱謝耶和華；我要在眾人中間讚美他；


\subsection{第110篇(大衛-基督是祭司-麦基洗德)：:耶和華對我主说的话和起的誓言 110:1-7}
\begin{table}[H]
\centering
\begin{tabular}{p{5.25cm}|p{12.2500cm}}\hline\hline
\multicolumn{1}{c|}{耶和華所說}&\multicolumn{1}{c}{我主所做}\\\hline
耶和華對我主說：你坐在我的右邊，等我使你仇敵作你的腳凳。
&耶和華必使你從錫安伸出能力的杖來；你要在你仇敵中掌權。
當你掌權的日子（行軍的日子），你的民要以聖潔的妝飾為衣（以聖潔為妝飾），甘心犧牲自己；你的民多如清晨的甘露（你少年時光耀如清晨的甘露）。\\\hline
耶和華起了誓決不後悔說：你是照著麥基洗德的等次永遠為祭司。
&在你右邊的主，當他發怒的日子必打傷列王。
他要在列邦中刑罰惡人，屍首就遍滿各處；他要在許多國中打破仇敵的頭。
他要喝路旁的河水，因此必擡起頭來。\\\hline
\end{tabular}
\end{table} 







\subsection{第111篇：在众民的会中赞美耶和华的作为 111:1-10}
\begin{table}[H]
\hspace{-1.5cm}
\centering
\begin{tabular}{p{3.75cm}|p{6.900cm}|p{4.00cm}|p{2.500cm}}\hline\hline
\multicolumn{1}{c|}{要讚美耶和華}&\multicolumn{1}{c|}{讚美所行的}&\multicolumn{1}{c|}{讚美他的訓詞和永約}&\multicolumn{1}{c}{遵行的是聰明人}\\\hline
你們要讚美耶和華！我要在正直人的大會中，並公會中，一心稱謝耶和華。
耶和華的作為本為大；凡喜愛的都必考察。
&他所行的是尊榮和威嚴；他的公義存到永遠。
他行了奇事，使人記念；耶和華有恩惠，有憐憫。
他賜糧食給敬畏他的人；他必永遠記念他的約。
他向百姓顯出大能的作為，把外邦的地賜給他們為業。他手所行的是誠實公平；
&他的訓詞都是確實的，
是永永遠遠堅定的，是按誠實正直設立的。
他向百姓施行救贖，命定他的約，直到永遠；他的名聖而可畏。
&敬畏耶和華是智慧的開端；凡遵行他命令的是聰明人。耶和華是永遠當讚美的！\\\hline
\end{tabular}
\end{table} 


\subsection{第112篇：敬畏遵守神命的福气 1-10}

\begin{table}[H]
\centering
\begin{tabular}{p{2.75cm}|p{3.500cm}|p{3.00cm}|p{5.00cm}|p{1.75cm}}\hline\hline
\multicolumn{1}{c|}{敬畏耶和華}&\multicolumn{1}{c|}{正直人}&\multicolumn{1}{c|}{施恩與人}&\multicolumn{1}{c|}{義人}&\multicolumn{1}{c}{惡人}\\\hline
你們要讚美耶和華！敬畏耶和華，甚喜愛他命令的，這人便為有福！
他的後裔在世必強盛；
&正直人的後代必要蒙福。
他家中有貨物，有錢財；他的公義存到永遠。
正直人在黑暗中，有光向他發現；他有恩惠，有憐憫，有公義。
&施恩與人、借貸與人的，這人事情順利；他被審判的時候要訴明自己的冤。
他永不動搖；
&義人被記念，直到永遠。
他必不怕凶惡的信息；他心堅定，倚靠耶和華。
他心確定，總不懼怕，直到他看見敵人遭報。
他施捨錢財，賙濟貧窮；他的仁義存到永遠。他的角必被高舉，大有榮耀。
&惡人看見便惱恨，必咬牙而消化；惡人的心願要歸滅絕。\\\hline
\end{tabular}
\end{table} 











\subsection{第113篇：讚美耶和华擡舉貧寒人 1-9}
\begin{table}[H]
\centering
\begin{tabular}{p{5.75cm}|p{11.00cm}}\hline\hline
\multicolumn{1}{c|}{你們要讚美耶和華 1-3}&\multicolumn{1}{c}{赞美耶和華的作为 4-9}\\\hline
\multirow{2}{5.9cm}{你們要讚美耶和華！耶和華的僕人哪，你們要讚美，讚美耶和華的名！耶和華的名是應當稱頌的，從今時直到永遠！
從日出之地到日落之處，耶和華的名是應當讚美的！}
&耶和華超乎萬民之上；他的榮耀高過諸天。
誰像耶和華─我們的神呢？他坐在至高之處，
自己謙卑，觀看天上地下的事。\\\cline{2-2}
&他從灰塵裡擡舉貧寒人，從糞堆中提拔窮乏人，
使他們與王子同坐，就是與本國的王子同坐。
他使不能生育的婦人安居家中，為多子的樂母。你們要讚美耶和華！\\\hline
\end{tabular}
\end{table} 








\subsection{第114篇：神对以色列家的救赎及受造之物的反应 114:1-8}
\begin{table}[H]
\centering
\begin{tabular}{p{5cm}|p{12.00cm}}\hline\hline
\multicolumn{1}{c|}{神对以色列家的救赎 1-2}&\multicolumn{1}{c}{受造之物的反应 3-8}\\\hline
以色列出了埃及，雅各家離開說異言之民；
&\multirow{3}{12.2cm}{\textcolor{red}{滄海}看見就奔逃；\textcolor{red}{約但河}也倒流。
\textcolor{red}{大山}踴躍，如公羊；\textcolor{red}{小山}跳舞，如羊羔。
滄海啊，你為何奔逃？約但哪，你為何倒流？
大山哪，你為何踴躍，如公羊？小山哪，你為何跳舞，如羊羔？
\textcolor{red}{大地}啊，你因見主的面，就是雅各神的面，便要震動。
他叫\textcolor{red}{磐石}變為水池，叫堅石變為泉源。}\\\cline{1-1}
那時，猶大為主的聖所，&\\\cline{1-1}
以色列為他所治理的國度。&\\\hline
\end{tabular}
\end{table} 



\subsection{第115篇：不要依靠外邦神，要依靠耶和华得福 115:1-18}
\subsubsection{不要依靠外邦偶像 1-8}
\begin{table}[H]
\centering
\begin{tabular}{p{6.75cm}|p{10.50cm}}\hline\hline
\multicolumn{1}{c|}{我們的神耶和華}&\multicolumn{1}{c}{外邦人和他們的偶像}\\\hline
耶和華啊，\textcolor{red}{榮耀}不要歸與我們，不要歸與我們；要因你的慈愛和誠實歸在你的名下！
&為何容外邦人說：他們的神在哪裡呢？\\\hline
然而，我們的\textcolor{red}{神在天上}，都\textcolor{red}{隨自己的意旨行事}。&他們的偶像是金的，銀的，是人手所造的，
有口卻不能言，有眼卻不能看，
有耳卻不能聽，有鼻卻不能聞，
有手卻不能摸，有腳卻不能走，有喉嚨也不能出聲。
造他的要和他一樣；凡靠他的也要如此。\\\hline
\end{tabular}
\end{table} 

\subsubsection{我们要依靠神 9-18}
\begin{table}[H]
\centering
\begin{tabular}{p{6.cm}|p{6.2500cm}|p{4.75cm}}\hline\hline
\multicolumn{1}{c|}{要依靠耶和华9-11}&\multicolumn{1}{c|}{耶和華要賜福的人12-15}&\multicolumn{1}{c}{蒙福后要讚美神16-18}\\\hline
以色列啊，你要倚靠耶和華！他是你的幫助和你的盾牌。
亞倫家啊，你們要倚靠耶和華！他是你們的幫助和你們的盾牌。
你們敬畏耶和華的，要倚靠耶和華！他是你們的幫助和你們的盾牌。
&耶和華向來眷念我們；他還要賜福給我們：要賜福給以色列的家，賜福給亞倫的家。
凡敬畏耶和華的，無論大小，主必賜福給他。
願耶和華叫你們和你們的子孫日見加增。
你們蒙了造天地之耶和華的福！
&天，是耶和華的天；地，他卻給了世人。
死人不能讚美耶和華；下到寂靜中的也都不能。
但我們要稱頌耶和華，從今時直到永遠。你們要讚美耶和華！\\\hline
\end{tabular}
\end{table} 





\subsection{第116篇：叙述神的拯救，及如何报答神的厚恩 1-18}
\subsubsection{讲述神如何垂听祷告和恳求 1-8}
\begin{table}[H]
\centering
\begin{tabular}{p{8.cm}|p{4.00cm}|p{4.75cm}}\hline\hline
\multicolumn{1}{c|}{求告耶和華的名 1-4}&\multicolumn{1}{c|}{救我脱离卑微 5-6}&\multicolumn{1}{c}{救我免了死亡流淚跌倒 7-8}\\\hline
我愛耶和華，因為他聽了我的聲音和我的懇求。
他既向我側耳，我一生要求告他。
死亡的繩索纏繞我；陰間的痛苦抓住我；我遭遇患難愁苦。
那時，我便求告耶和華的名，說：耶和華啊，求你救我的靈魂！
&耶和華有恩惠，有公義；我們的神以憐憫為懷。
耶和華保護愚人；我落到卑微的地步，他救了我。
&我的心哪！你要仍歸安樂，因為耶和華用厚恩待你。
主啊，你救我的命免了死亡，救我的眼免了流淚，救我的腳免了跌倒。\\\hline
\end{tabular}
\end{table} 


\subsubsection{我要如何回应神的厚恩 9-18}
我要在耶和華面前行活人之路。
我因信，所以如此說話；我受了極大的困苦。
我曾急促地說：人都是說謊的！

我拿什麼報答耶和華向我所賜的一切厚恩？
我要舉起救恩的杯，稱揚耶和華的名。
我要在他眾民面前向耶和華還我的願。
在耶和華眼中，看聖民之死極為寶貴。
耶和華啊，我真是你的僕人；我是你的僕人，是你婢女的兒子。你已經解開我的綁索。
我要以感謝為祭獻給你，又要求告耶和華的名。
我要在他眾民面前，在耶和華殿的院內，在耶路撒冷當中，向耶和華還我的願。你們要讚美耶和華！



\subsection{第117篇：萬國萬民都要赞美神的慈愛和誠實 1-2}
萬國啊，你們都當讚美耶和華！萬民哪，你們都當頌讚他！
因為他向我們大施慈愛；耶和華的誠實存到永遠。你們要讚美耶和華！

\subsection{第118篇(基督的被弃高升):因神拯救人的计划，呼吁所有人都要赞美神 1-29}
\subsubsection{要稱謝耶和華的善和永遠長存的慈愛 1-4}
你們要稱謝耶和華，因他本為善；他的慈愛永遠長存！
願\textcolor{red}{以色列}說：他的慈愛永遠長存！
願\textcolor{red}{亞倫的家}說：他的慈愛永遠長存！
願\textcolor{red}{敬畏耶和華的}說：他的慈愛永遠長存！

\subsubsection{述說神拯救我的见证 5-21}
我在急難中求告耶和華，他就應允我，把我安置在寬闊之地。
有耶和華幫助我，我必不懼怕，人能把我怎麼樣呢？
在那幫助我的人中，有耶和華幫助我，所以我要看見那恨我的人遭報。
投靠耶和華，強似倚賴人；
投靠耶和華，強似倚賴王子。

萬民圍繞我，我靠耶和華的名必剿滅他們。
他們環繞我，圍困我，我靠耶和華的名必剿滅他們。
他們如同蜂子圍繞我，好像燒荊棘的火，必被熄滅；我靠耶和華的名，必剿滅他們。
你推我，要叫我跌倒，但耶和華幫助了我。
耶和華是我的力量，是我的詩歌；他也成了我的拯救。

在義人的帳棚裡，有歡呼拯救的聲音；耶和華的右手施展大能。
耶和華的右手高舉；耶和華的右手施展大能。
我必不致死，仍要存活，並要傳揚耶和華的作為。
耶和華雖嚴嚴地懲治我，卻未曾將我交於死亡。
給我敞開義門；我要進去稱謝耶和華！
這是耶和華的門；義人要進去！
我要稱謝你，因為你已經應允我，又成了我的拯救！

\subsubsection{述說神拯救人类的计划，呼吁大家称谢耶和华 22-29}
匠人所棄的石頭已成了房角的頭塊石頭。
這是耶和華所做的，在我們眼中看為希奇。
這是耶和華所定的日子，我們在其中要高興歡喜！
耶和華啊，求你拯救！耶和華啊，求你使我們亨通！
奉耶和華名來的是應當稱頌的！我們從耶和華的殿中為你們祝福！
耶和華是神；他光照了我們。理當用繩索把祭牲拴住，牽到壇角那裡。
你是我的神，我要稱謝你！你是我的神，我要尊崇你！
你們要稱謝耶和華，因他本為善；他的慈愛永遠長存！

\subsection{第119篇：神的法度/律例/话/典章 119:1-176}
%\subsubsection{Aleph \<'a+s:dey>：遵行耶和華律法的(有福) 1-8}
\subsubsection{Aleph \texorpdfstring{\ez{א}}{\unichar{"05D0}}：遵行耶和華律法的(\textcolor{red}{有福}) 1-8}
\begin{itemize}
\itemsep-.45em
\item 有福
\begin{itemize}
\itemsep-.45em
\item 行為完全、遵行耶和華律法的有福
\item 遵守他的法度、一心尋求他,不做非義的事，但遵行他的道，有福
\end{itemize}
\item 神將訓詞吩咐我們的\underline{目的}，為要我們殷勤遵守
\item 遵守神律法的结局
\begin{itemize}
\itemsep-.45em
\item 因\underline{遵守}你的律例，願我人生道路就\underline{堅定}
\item \underline{看重}你的一切命令，就\underline{不羞愧}
\item \underline{學了}你公義的判語，就要以正直的心\underline{稱謝}你
\item 必\underline{守}你的律例；求你總\underline{不丟棄}我
\end{itemize}
\end{itemize}
行為完全、遵行耶和華律法的，這人便為有福！
遵守他的法度、一心尋求他的，這人便為有福！
這人不做非義的事，但遵行他的道。
耶和華啊，你曾將你的訓詞吩咐我們，為要我們殷勤遵守。
但願我行事堅定，得以遵守你的律例。
我看重你的一切命令，就不至於羞愧。
我學了你公義的判語，就要以正直的心稱謝你。
我必守你的律例；求你總不要丟棄我！


%\subsubsection{Beth \<b*am*EHh--> ：少年人用守神的話（怎么，什么）潔淨行為 9-16}
\subsubsection{Beth \texorpdfstring{\ez{ב}}{\unichar{"05D1}}:少年人用守神的話（\textcolor{red}{怎么，什么}）潔淨行為 9-16}
\begin{itemize}
\itemsep-.45em
\item 为什么洁净少年人的心
\begin{itemize}
\itemsep-.45em
\item 少年人用遵守神的话\underline{潔淨行為}
\item 一心尋求你；求你叫我\underline{不偏離}你的命令
\item 將你的話藏在心裡，\underline{免得我得罪你}
\end{itemize}
\item 如何洁净少年人的心
\begin{itemize}
\itemsep-.45em
\item 求\underline{你教訓我}你的律例,耶和華是應當稱頌的
\item 用嘴唇\underline{傳揚}你口中的一切典章
\item \underline{喜悅}你的法度，如同喜悅一切的財物
\item \underline{默想}你的訓詞，看重你的道路
\item 在你的律例中\underline{自樂}；我\underline{不忘記}你的話
\end{itemize}
\end{itemize}
少年人用什麼潔淨他的行為呢？是要遵行你的話！
我一心尋求了你；求你不要叫我偏離你的命令。
我將你的話藏在心裡，免得我得罪你。
耶和華啊，你是應當稱頌的！求你將你的律例教訓我！
我用嘴唇傳揚你口中的一切典章。
我喜悅你的法度，如同喜悅一切的財物。
我要默想你的訓詞，看重你的道路。
我要在你的律例中自樂；我不忘記你的話。


%\subsubsection{Gimel \<g*:mol>：求你用\protect{\textcolor{red}{厚恩待/恶待}}我，使我存活遵守你的話 17-24}
\subsubsection{Gimel \texorpdfstring{\ez{ג}}{\unichar{"05D2}}:{\textcolor{red}{厚恩待/恶待}}我，使我存活遵守你的話 17-24}
\begin{itemize}
\itemsep-.45em
\item 向神的祈求 17-20
\begin{itemize}
\itemsep-.45em
\item 求你用厚恩待你的僕人，使我存活，我就遵守你的話
\item 求你開我的眼睛，使我看出你律法中的奇妙
\item 求你不要向我隱瞞你的命令,我是在地上作寄居的.我時常切慕你的典章，甚至心碎
\end{itemize}
\item 在首领妄论我和受羞辱时我遵守你的法度 21-24
\begin{itemize}
\itemsep-.45em
\item 受咒詛、偏離你命令的驕傲人，你已經責備他們
\item 求你除掉我所受的羞辱和藐視，因我遵守你的法度
\item 雖有首領坐著妄論我，你僕人卻思想你的律例,你的法度是我所喜樂的，是我的謀士
\end{itemize}
\end{itemize}
求你用厚恩待你的僕人，使我存活，我就遵守你的話。
求你開我的眼睛，使我看出你律法中的奇妙。
我是在地上作寄居的；求你不要向我隱瞞你的命令！
我時常切慕你的典章，甚至心碎。
受咒詛、偏離你命令的驕傲人，你已經責備他們。
求你除掉我所受的羞辱和藐視，因我遵守你的法度。
雖有首領坐著妄論我，你僕人卻思想你的律例。
你的法度是我所喜樂的，是我的謀士。



%\subsubsection{Daleth \<d*Ab:qAh>： \textcolor{red}{Cleave(劈开，开辟）}在艰难和愁苦揀選,忠信，忠心 25-32}
\subsubsection{Daleth \texorpdfstring{\ez{ד}}{\unichar{"05D3}}:\textcolor{red}{Cleave(劈开，开辟）}在艰难和愁苦揀選,忠信，忠心 25-32}
\begin{itemize}
\itemsep-.45em
\item 我的性命幾乎歸於塵土；求你照你的話將我救活！ 
\begin{itemize}
\itemsep-.45em
\item  我述說我所行的，你應允了我；求你將你的律例教訓我！ 
\item  求你使我明白你的訓詞，我就思想你的奇事。
\end{itemize}
\item  我的心因愁苦而消化；求你照你的話使我堅立！
\begin{itemize}
\itemsep-.45em 
\item  求你使我離開奸詐的道，開恩將你的律法賜給我！ 
\item  我揀選了忠信的道，將你的典章擺在我面前。 
\item  我持守你的法度；耶和華啊，求你不要叫我羞愧！ 
\item  你開廣我心的時候，我就往你命令的道上直奔
\end{itemize}
\end{itemize}
我的性命幾乎歸於塵土；求你照你的話將我救活！
我述說我所行的，你應允了我；求你將你的律例教訓我！
求你使我明白你的訓詞，我就思想你的奇事。
我的心因愁苦而消化；求你照你的話使我堅立！
求你使我離開奸詐的道，開恩將你的律法賜給我！
我揀選了忠信的道，將你的典章擺在我面前。
我持守你的法度；耶和華啊，求你不要叫我羞愧！
你開廣我心的時候，我就往你命令的道上直奔。



%\subsubsection{Hey \<hworeniy>： \textcolor{red}{Teach me 教导}我,我必遵守到底 33-40}
\subsubsection{Hey \texorpdfstring{\ez{ה}}{\unichar{"05D4}}:\textcolor{red}{Teach me 教导}我,我必遵守到底 33-40}
耶和華啊，求你將你的律例指教我，我必遵守到底！
求你賜我悟性，我便遵守你的律法，且要一心遵守。
求你叫我遵行你的命令，因為這是我所喜樂的。
求你使我的心趨向你的法度，不趨向非義之財。
求你叫我轉眼不看虛假，又叫我在你的道中生活。
你向敬畏你的人所應許的話，求你向僕人堅定！
求你使我所怕的羞辱遠離我，因你的典章本為美。
我羨慕你的訓詞；求你使我在你的公義上生活！




%\subsubsection{Waw \<wiybo'uniy>：\textcolor{red}{Let come to me 愿}慈愛救恩\textcolor{red}{临到我} 41-48}
\subsubsection{Waw \texorpdfstring{\ez{ו}}{\unichar{"05D5}}:\textcolor{red}{Let come to me 愿}慈愛救恩\textcolor{red}{临到我} 41-48}
耶和華啊，願你照你的話，使你的慈愛，就是你的救恩，臨到我身上，
我就有話回答那羞辱我的，因我倚靠你的話。
求你叫真理的話總不離開我口，因我仰望你的典章。
我要常守你的律法，直到永永遠遠。
我要自由而行（或譯：我要行在寬闊之地），因我素來考究你的訓詞。
我也要在君王面前論說你的法度，並不至於羞愧。
我要在你的命令中自樂；這命令素來是我所愛的。
我又要遵行（原文是舉手）你的命令，這命令素來是我所愛的；我也要思想你的律例。


%\subsubsection{Zqyin  \<z:kor-->：\textcolor{red}{Remember記念}向你僕人的應許叫我有盼望 49-56}
\subsubsection{Zqyin \texorpdfstring{\ez{ז}}{\unichar{"05D6}}:\textcolor{red}{Remember記念}向你僕人的應許叫我有盼望 49-56}
求你記念向你僕人所應許的話，叫我有盼望。
這話將我救活了；我在患難中，因此得安慰。
驕傲的人甚侮慢我，我卻未曾偏離你的律法。
耶和華啊，我記念你從古以來的典章，就得了安慰。
我見惡人離棄你的律法，就怒氣發作，猶如火燒。
我在世寄居，素來以你的律例為詩歌。
耶和華啊，我夜間記念你的名，遵守你的律法。
我所以如此，是因我守你的訓詞。



%\subsubsection{Heth  \<.hEl:qiy>：耶和華是\textcolor{red}{我的福分my portion} 57-64}
\subsubsection{Heth \texorpdfstring{\ez{ח}}{\unichar{"05D7}}:耶和華是\textcolor{red}{我的福分my portion} 57-64}
耶和華是我的福分；我曾說，我要遵守你的言語。 
我一心求過你的恩；願你照你的話憐憫我！ 
Psa 119:59  我思想我所行的道，就轉步歸向你的法度。 
Psa 119:60  我急忙遵守你的命令，並不遲延。 
Psa 119:61  惡人的繩索纏繞我，我卻沒有忘記你的律法。 
Psa 119:62  我因你公義的典章，半夜必起來稱謝你。 
Psa 119:63  凡敬畏你、守你訓詞的人，我都與他作伴。 
Psa 119:64  耶和華啊，你的慈愛遍滿大地；求你將你的律例教訓我


%\subsubsection{Teth  \<.twob>：你向來是照你的話\textcolor{red}{well善待}僕人 65-72}
\subsubsection{Tet \texorpdfstring{\ez{ט}}{\unichar{"05D8}}:你向來是照你的話\textcolor{red}{well善待}僕人 65-72}
\begin{itemize}
\itemsep-.45em
\item 耶和華啊，你向來是照你的話善待僕人。 
Psa 119:66  求你將精明和知識賜給我，因我信了你的命令。 
Psa 119:67  我未受苦以先走迷了路，現在卻遵守你的話。 
Psa 119:68  你本為善，所行的也善；求你將你的律例教訓我！ 
Psa 119:69  驕傲人編造謊言攻擊我，我卻要一心守你的訓詞。 
Psa 119:70  他們心蒙脂油，我卻喜愛你的律法。 
Psa 119:71  我受苦是與我有益，為要使我學習你的律例。 
Psa 119:72  你口中的訓言（律法）與我有益，勝於千萬的金銀
\end{itemize}

%\subsubsection{Yodh \<yAdEyKA>：\textcolor{red}{Your hands你的手}建立我，求賜悟性學習你的命令 73-80}
\subsubsection{Yodh \texorpdfstring{\ez{י}}{\unichar{"05D9}}:\textcolor{red}{Your hands你的手}建立我，求賜悟性學習你的命令 73-80}
你的手製造我，建立我；求你賜我悟性，可以學習你的命令！ 
敬畏你的人見我就要歡喜，因我仰望你的話。 
耶和華啊，我知道你的判語是公義的；你使我受苦是以誠實待我。 
求你照著應許僕人的話，以慈愛安慰我。 
願你的慈悲臨到我，使我存活，因你的律法是我所喜愛的。 
願驕傲人蒙羞，因為他們無理地傾覆我；但我要思想你的訓詞。 
願敬畏你的人歸向我，他們就知道你的法度。 
願我的心在你的律例上完全，使我不至蒙羞



%\subsubsection{Caph \<kAl:tiAh>：我心\textcolor{red}{渴想faints}你的救恩，仰望你的應許 81-88}
\subsubsection{Caph \texorpdfstring{\ez{כ}}{\unichar{"05DB}}:我心\textcolor{red}{渴想faints}你的救恩，仰望你的應許 81-88}
我心渴想你的救恩，仰望你的應許。
我因盼望你的應許眼睛失明，說：你何時安慰我？
我好像煙薰的皮袋，卻不忘記你的律例。
你僕人的年日有多少呢？你幾時向逼迫我的人施行審判呢？
不從你律法的驕傲人為我掘了坑。
你的命令盡都誠實；他們無理地逼迫我，求你幫助我！
他們幾乎把我從世上滅絕，但我沒有離棄你的訓詞。
求你照你的慈愛將我救活，我就遵守你口中的法度。


\subsubsection{Lamed \texorpdfstring{\ez{ל}}{\unichar{"05DC}}:你的話安定在天，直到\textcolor{red}{永遠 Forever} 89-96}

%\subsubsection{Lamed \<l:`wolAm>：你的話安定在天，直到\textcolor{red}{永遠 Forever} 89-96}
%\subsubsection{Lamed \texorpdfstring{\ez{ל}}{\unichar{"05DC}}:你的話安定在天，直到\textcolor{red}{永遠 Forever} 89-96}
耶和華啊，你的話安定在天，直到永遠。
你的誠實存到萬代；你堅定了地，地就長存。
天地照你的安排存到今日；萬物都是你的僕役。
我若不是喜愛你的律法，早就在苦難中滅絕了！
我永不忘記你的訓詞，因你用這訓詞將我救活了。
我是屬你的，求你救我，因我尋求了你的訓詞。
惡人等待我，要滅絕我，我卻要揣摩你的法度。
我看萬事盡都有限，惟有你的命令極其寬廣。


%\subsubsection{Mem \<mAh-->：我\textcolor{red}{何等}愛慕你的律法，終日不住地思想 97-104}
\subsubsection{Mem \texorpdfstring{\ez{ם}}{\unichar{"05DD}}:我\textcolor{red}{何等}愛慕你的律法，終日不住地思想 97-104}
我何等愛慕你的律法，終日不住地思想。
你的命令常存在我心裡，使我比仇敵有智慧。
我比我的師傅更通達，因我思想你的法度。
我比年老的更明白，因我守了你的訓詞。
我禁止我腳走一切的邪路，為要遵守你的話。
我沒有偏離你的典章，因為你教訓了我。
你的言語在我上膛何等甘美，在我口中比蜜更甜！
我藉著你的訓詞得以明白，所以我\textcolor{blue}{恨一切的假道}。


%\subsubsection{Nun \<ned-->：你的話是我腳前的\textcolor{red}{燈}，是我路上的光 105-112}
\subsubsection{Nun \texorpdfstring{\ez{נ}}{\unichar{"05E0}}:你的話是我腳前的\textcolor{red}{燈}，是我路上的光 105-112}
你的話是我腳前的燈，是我路上的光。
你公義的典章，我曾起誓遵守，我必按誓而行。
我甚是受苦；耶和華啊，求你照你的話將我救活！
耶和華啊，求你悅納我口中的讚美為供物，又將你的典章教訓我！
我的性命常在危險之中，我卻不忘記你的律法。
惡人為我設下網羅，我卻沒有偏離你的訓詞。
我以你的法度為永遠的產業，因這是我心中所喜愛的。
我的心專向你的律例，永遠遵行，一直到底。


%\subsubsection{Samech \<se`a:piyM>：心懷二意\textcolor{red}{[vain] thoughts}的人為我所恨；你的律法為我所愛 113-120}
\subsubsection{Samech \texorpdfstring{\ez{ס}}{\unichar{"05E1}}:心懷二意\textcolor{red}{[vain] thoughts}的人為我所恨；你的律法為我所愛 113-120}
心懷二意的人為我所恨；但你的律法為我所愛。
你是我藏身之處，又是我的盾牌；我甚仰望你的話語。
作惡的人哪，你們離開我吧！我好遵守我神的命令。
求你照你的話扶持我，使我存活，也不叫我因失望而害羞。
求你扶持我，我便得救，時常看重你的律例。
凡偏離你律例的人，你都輕棄他們，因為他們的詭詐必歸虛空。
凡地上的惡人，你除掉他，好像除掉渣滓；因此我愛你的法度。
我因懼怕你，肉就發抖；我也怕你的判語。


%\subsubsection{Ain \<`A,siytiy>：\textcolor{red}{我行過I have done}公平和公義，不要撇下我給欺壓我的人 121-128 }
\subsubsection{Ain \texorpdfstring{\ez{ע}}{\unichar{"05E2}}:\textcolor{red}{我行過I have done}公平和公義，不要撇下我給欺壓我的人 121-128 }
我行過公平和公義，求你不要撇下我給欺壓我的人！
求你為僕人作保，使我得好處，不容驕傲人欺壓我！
我因盼望你的救恩和你公義的話眼睛失明。
求你照你的慈愛待僕人，將你的律例教訓我。
我是你的僕人，求你賜我悟性，使我得知你的法度。
這是耶和華降罰的時候，因人廢了你的律法。
所以，我愛你的命令勝於金子，更勝於精金。
你一切的訓詞，在萬事上我都以為正直；我卻\textcolor{blue}{恨惡一切假道}。

%\subsubsection{Pe \<p*:lA'wot> ：你的法度\textcolor{red}{[are]wonderful 是奇妙}我一心謹守 129-136}
\subsubsection{Pe \texorpdfstring{\ez{פ}}{\unichar{"05E4}}:你的法度\textcolor{red}{[are]wonderful 是奇妙}我一心謹守 129-136}
你的法度奇妙，所以我一心謹守。
你的言語一解開就發出亮光，使愚人通達。
我張口而氣喘，因我切慕你的命令。
求你轉向我，憐憫我，好像你素常待那些愛你名的人。
求你用你的話使我腳步穩當，不許什麼罪孽轄制我。
求你救我脫離人的欺壓，我要遵守你的訓詞。
求你用臉光照僕人，又將你的律例教訓我。
我的眼淚下流成河，因為他們不守你的律法。


%\subsubsection{Tzaddi \<.sad*iyq>：你是\textcolor{red}{公義}的；你的判語也是正直的 137-144}
\subsubsection{Tzaddi \texorpdfstring{\ez{צ}}{\unichar{"05E6}}:你是\textcolor{red}{公義}的；你的判語也是正直的 137-144}
耶和華啊，你是公義的；你的判語也是正直的！
你所命定的法度是憑公義和至誠。
我心焦急，如同火燒，因我敵人忘記你的言語。
你的話極其精煉，所以你的僕人喜愛。
我微小，被人藐視，卻不忘記你的訓詞。
你的公義永遠長存；你的律法盡都真實。
我遭遇患難愁苦，你的命令卻是我所喜愛的。
你的法度永遠是公義的；求你賜我悟性，我就活了。

%\subsubsection{Koph \<qArA'tiy>：我一心呼籲\textcolor{red}{I cried 我哭喊}你；求應允我，我必謹守律例 145-152}
\subsubsection{Koph \texorpdfstring{\ez{ק}}{\unichar{"05E7}}:我一心呼籲\textcolor{red}{I cried 我哭喊}你；求應允我，我必謹守律例 145-152}
耶和華啊，我一心呼籲你；求你應允我，我必謹守你的律例！
我向你呼籲，求你救我！我要遵守你的法度。
我趁天未亮呼求；我仰望了你的言語。
我趁夜更未換將眼睜開，為要思想你的話語。
求你照你的慈愛聽我的聲音；耶和華啊，求你照你的典章將我救活！
追求奸惡的人臨近了；他們遠離你的律法。
耶和華啊，你與我相近；你一切的命令盡都真實！
我因學你的法度，久已知道是你永遠立定的。

%\subsubsection{Resh  \<d:'eh-->：求你看顧搭救我，因我不忘記(\textcolor{red}{Consider思想})你的律法  153-160}
\subsubsection{Resh \texorpdfstring{\ez{ר}}{\unichar{"05E8}}:求你看顧搭救我，因我不忘記(\textcolor{red}{Consider思想})你的律法  153-160}
求你看顧我的苦難，搭救我，因我不忘記你的律法。
求你為我辨屈，救贖我，照你的話將我救活。
救恩遠離惡人，因為他們不尋求你的律例。
耶和華啊，你的慈悲本為大；求你照你的典章將我救活。
逼迫我的，抵擋我的，很多，我卻沒有偏離你的法度。
我看見奸惡的人就甚憎惡，因為他們不遵守你的話。
你看我怎樣愛你的訓詞！耶和華啊，求你照你的慈愛將我救活！
你話的總綱是真實；你一切公義的典章是永遠長存。


%\subsubsection{Schin \<,sAdiyM>：首領(\textcolor{red}{Princes 王子})無故地逼迫我，但我的心畏懼你的言語  161-168}
\subsubsection{Schin \texorpdfstring{\ez{ש}}{\unichar{"05E9}}:首領(\textcolor{red}{Princes 王子})無故地逼迫我，但我的心畏懼你的言語  161-168}
首領無故地逼迫我，但我的心畏懼你的言語。
我喜愛你的話，好像人得了許多擄物。
謊話是我所恨惡所憎嫌的；惟你的律法是我所愛的。
我因你公義的典章一天七次讚美你。
愛你律法的人有大平安，什麼都不能使他們絆腳。
耶和華啊，我仰望了你的救恩，遵行了你的命令。
我心裡守了你的法度；這法度我甚喜愛。
我遵守了你的訓詞和法度，因我一切所行的都在你面前。


%\subsubsection{Tau \<t*iq:dab>：願我的呼籲\textcolor{red}{達到(Let come near)}你面前，照你的話賜我悟性 169-176}
\subsubsection{Tau  \texorpdfstring{\ez{ת}}{\unichar{"05EA}}:願我的呼籲\textcolor{red}{達到(Let come near)}你面前，照你的話賜我悟性 169-176}
耶和華啊，願我的呼籲達到你面前，照你的話賜我悟性。
願我的懇求達到你面前，照你的話搭救我。
願我的嘴發出讚美的話，因為你將律例教訓我。
願我的舌頭歌唱你的話，因你一切的命令盡都公義。
願你用手幫助我，因我揀選了你的訓詞。
耶和華啊，我切慕你的救恩！你的律法也是我所喜愛的。
願我的性命存活，得以讚美你！願你的典章幫助我！
我如亡羊走迷了路，求你尋找僕人，因我不忘記你的命令。





























